
\section{Comparación de las propiedades de los morfismos}
\label{XII.3}

\begin{proposition}
\label{XII.3.1}
Sean $f\colon X\to Y$ un morfismo de $\CC$-esquemas localmente de tipo
finito y $f^\an\colon X^\an\to Y^\an$ el morfismo deducido de~$f$ sobre los
espacios analíticos asociados. Sea~$P$ la propiedad de ser
\begin{enumerate}
\item[(i)] plano
\item[(ii)] neto (\ie no ramificado)
\item[(iii)] étale
\item[(iv)] liso
\item[(v)] normal
\item[(vi)]
\marginpar{321}
reducido
\item[(vii)] inyectivo
\item[(viii)] separado
\item[(ix)] un isomorfismo
\item[(x)] un monomorfismo
\item[(xi)] una inmersión abierta.
\end{enumerate}
Entonces, para que~$f$ posea la propiedad~$P$, es necesario y suficiente que
también la posea~$f^\an$.
\end{proposition}

Denotemos por $\varphi\colon X^\an\to X$ y $\psi\colon Y^\an\to Y$ los
morfismos canónicos. Sean~$x$ un punto de~$X^\an$ e $y=f^\an(x)$. Los
morfismos $\cal{O}_{Y^\an,y}\to{\mathcal O}_{X^\an,x}$ y
$\cal{O}_{Y,\psi(y)}\to\cal{O}_{X,\varphi(x)}$,
\ifthenelse{\boolean{orig}}
{inducidos por~$f$ y~$f^\an$, respectivamente,}
{inducidos por~$f^\an$ y~$f$, respectivamente,}%
dan el mismo morfismo al pasar a las completaciones~\eqref{XII.1.1}. Según
\cite[ch\ptbl 3 \S 5 prop\ptbl 4]{XII.2} (\resp EGA~IV~17.4.4), equivale,
pues, decir que~$f^\an$ verifica la propiedad~(i) (\resp (ii)) o decir que~$f$
verifica~(i) (\resp (ii)) en cada punto cerrado de~$X$. Como el conjunto de
los puntos de~$X$ en los que se verifica~(i) (\resp (ii)) es abierto
(EGA~IV~11.1.1 e I~3.3), esto demuestra~(i) y~(ii) y, por tanto,
también~(iii).

\ifthenelse{\boolean{orig}}{}
{\enlargethispage{.5cm}}%
Sea~$P$ la propiedad~(iv) (\resp (v), \resp (vi)). Teniendo en cuenta
\Ref{XII.2.1}
\ifthenelse{\boolean{orig}}
{((v), (vi), (vii)),}
{((iv), (vi), (vii)),}%
~equivale decir que las fibras geométricas de~$f^\an$ en los distintos
puntos~$y$ de~$Y^\an$ son regulares (\resp normales, \resp reducidas) o que
lo mismo ocurre con las fibras geométricas de~$f$ en los distintos puntos
cerrados~$\psi(y)$ de~$Y$. Los casos~(iv) (\resp (v), \resp (vi)) resultan
entonces de~(i) y del hecho de que el conjunto de los puntos de~$Y$ en los
que las fibras geométricas de~$f$ son regulares es abierto
(EGA~IV~12.1.7).

(vii). Si~$f$ es inyectivo, también lo es~$f^\an$. Recíprocamente,
supongamos que~$f^\an$ es inyectivo y demostremos que~$f$ también lo es.
Podemos suponer
\marginpar{322}
que~$f$ es de tipo finito. Como el morfismo~$f^\an$ es inyectivo, las fibras
de~$f$ en los puntos cerrados de~$Y$ son radiciales; como el conjunto de los
puntos de~$Y$ cuya fibra es radicial es localmente constructible
(EGA~IV~9.6.1) y como~$Y$ es un esquema de Jacobson, todas las fibras de~$f$
son radiciales y, por tanto, es inyectivo.

(viii). Sean $\Delta\colon X\to X\times_YX$ y
$\Delta^\an\colon X^\an\to X^\an\times_{Y^\an}X^\an$ las inmersiones
diagonales, y sea
$\Theta\colon X^\an\times_{Y^\an}X^\an\to X\times_YX$ el morfismo canónico.
En virtud de~\Ref{XII.2.3}, equivale decir que~$\Delta(X)$ es cerrado en
$X\times_YX$ o que~$\Delta^\an(X^\an)$ es cerrado en
$X^\an\times_{Y^\an}X^\an$.

Como una inmersión abierta no es sino un morfismo étale inyectivo
(EGA~IV~17.9.1 y~\cite[\no13 \S 1]{XII.4}), (xi) resulta de~(iii) y~(vii).
Como un isomorfismo es lo mismo que una inmersión abierta sobreyectiva, (ix)
resulta de~(xi) y de (\Ref{XII.3.2}~(i)) más abajo. Decir que~$f$ es un
monomorfismo equivale a decir que
\ifthenelse{\boolean{orig}}
{la morfismo}
{el morfismo}%
~diagonal $\Delta\colon X\to X\times_YX$ es un isomorfismo; por tanto, (x)
resulta de~(ix).

\begin{proposition}
\label{XII.3.2}
Sean~$X$ e~$Y$ dos $\CC$-esquemas localmente de tipo finito,
$f\colon X\to Y$ un morfismo de tipo finito y
$f^\an\colon X^\an\to Y^\an$ el morfismo deducido de~$f$ sobre los espacios
analíticos asociados. Sea~$P$ la propiedad de ser
\begin{enumerate}
\item[(i)] sobreyectivo
\item[(ii)] dominante
\item[(iii)] una inmersión cerrada
\item[(iv)] una inmersión
\item[(v)] propio\footnote{Diremos que un morfismo de espacios analíticos es
propio si lo es en el sentido de~\cite[ch\ptbl 1 \S 10 \no1]{XII.1} y si es
separado.}
\item[(vi)]
\marginpar{323}
finito.
\end{enumerate}
Entonces, para que~$f$ posea la propiedad~$P$, es necesario y suficiente que
también la posea~$f^\an$.
\end{proposition}

Sean $\varphi\colon X^\an\to X$ y $\psi\colon Y^\an\to Y$ los morfismos
canónicos.

(i). Si~$f$ es sobreyectivo, para todo punto~$y$ de~$Y^\an$,
$f^{-1}(\psi(y))$ es un subconjunto cerrado no vacío de~$X$; contiene, pues,
al menos un punto cerrado, lo cual demuestra que~$f^\an$ es sobreyectivo.
Recíprocamente, si~$f^\an$ es sobreyectivo, $f(X)$ es un subconjunto
localmente constructible de~$Y$ (EGA~IV~1.8.4) que contiene todos los puntos
cerrados de~$Y$; por tanto, $f(X)=Y$.

(ii) resulta de~\Ref{XII.2.2}.

(iii). Si~$f$ es una inmersión cerrada, también lo es~$f^\an$ según
\Ref{XII.1.1}~a). Recíprocamente, si~$f^\an$ es una inmersión cerrada, también
lo es~$f$ según~\Ref{XII.3.1}~(x) y~\Ref{XII.3.2}~(v), pues esto equivale a
decir que~$f$ es un monomorfismo propio (EGA~IV~8.11.5).

(iv). Es claro que, si~$f$ es una inmersión, también lo es~$f^\an$.
Recíprocamente, supongamos que~$f^\an$ es una inmersión, y sean~$T$ la imagen
de~$X$ en~$Y$ y~$\overline{T}$ la clausura esquemática de~$f$. Se tiene una
factorización de~$f$,
$$
X \lto{i} \overline{T} \lto{j}
\ifthenelse{\boolean{orig}}
{J}
{Y}%
\quoi,
$$
donde~$j$ es una inmersión cerrada e~$i$ el morfismo canónico; de ella se
deduce la factorización siguiente de~$f^\an$:
$$
X^\an \lto{i^\an} \overline{T}^{\mkern1mu \an} \lto{j^\an} Y^\an\ .
$$
Como
\marginpar{324}
$T=f(X)$ es un subconjunto localmente constructible de~$Y$ (EGA~IV~1.8.4),
se tiene, según~\Ref{XII.2.2},
$\overline{T}^{\mkern1mu \an}=\overline{f^\an(X^\an)}$. De ello resulta que
$i^\an(X^\an)$ es un abierto de~$\overline{T}^{\mkern1mu \an}$ y, por tanto,
que~$i(X)$ es un abierto de~$\overline{T}$. Consideremos la factorización
canónica de~$i$:
$$
X \lto{i_1} i(X) \lto{i_2} \overline{T}.
$$
El morfismo~$i_1^\an$ es un monomorfismo propio; por tanto, también lo es
$i_1$ según~\Ref{XII.3.2}~(v) y~\Ref{XII.3.1}~(x). Esto demuestra que~$i_1$
y, por consiguiente, también~$f$ es una inmersión.

(v). Supongamos que~$f$ es propio y demostremos que~$f^\an$ también lo es.
Como el hecho de que~$f^\an$ sea propio es local sobre~$Y^\an$, podemos
suponer que~$Y$ es afín. Según el lema de Chow (EGA~II~5.6.1), puede
encontrarse un $Y$-esquema proyectivo~$X'$ y un morfismo proyectivo
sobreyectivo
$$
g\colon X'\to X.
$$
El morfismo $(fg)^\an=f^\an g^\an$ es proyectivo y, por tanto, propio;
$g^\an$ es sobreyectivo, y de~\cite[ch\ptbl 1 \S 10]{XII.1} resulta que
$f^\an$ es propio.

Recíprocamente, supongamos que~$f^\an$ es propio y demostremos que~$f$
también lo es. Según~\Ref{XII.3.1}~(viii), $f$ es separado. Falta probar que
$f$ es universalmente cerrado, e incluso basta demostrar que~$f$ es cerrado;
en efecto, para todo $Y$-esquema~$Y'$ localmente de tipo finito, el morfismo
$$
f_{(Y')}=h\colon X\times_YY'\to Y'
$$
también será cerrado, pues~$h^\an$ es propio. Sea~$T$ un subconjunto cerrado
de~$X$; $f(T)$ es un conjunto localmente constructible y se tiene
$$
f^\an(\varphi^{-1}(T))=\psi^{-1}(f(T)).
$$
Como
\marginpar{325}
$f^\an$ es propio, $\psi^{-1}(f(T))$ es un subconjunto cerrado de~$Y^\an$;
de~\Ref{XII.2.2} resulta, pues, que
$$
\psi^{-1}(\overline{f(T)})=\psi^{-1}(f(T))\quoi.
$$
Esto implica que $f(T)=\overline{f(T)}$, \ie que~$f$ es
\ifthenelse{\boolean{orig}}
{cerrada}
{cerrado}%
~y, por tanto, que~$f$ es propio.

(vi). Decir que un morfismo es finito equivale a decir que es propio de
fibras finitas (EGA~III~4.4.2 y~\cite[\no19 \S 5]{XII.4}). Como el conjunto
de los puntos en los que las fibras de~$f$ son finitas es localmente
constructible (EGA~IV~9.7.9), las fibras de~$f$ son finitas
\ifthenelse{\boolean{orig}}
{si y solo si}
{si y solo si}%
~lo son las fibras de~$f^\an$; por tanto, (vi) resulta de~(v).

\begin{remarque}
\label{XII.3.3}
\begin{enumerate}
\item[a)] sea $f\colon X\to Y$ un morfismo de $\CC$-esquemas localmente de
tipo finito. El hecho de que~$f^\an$ sea un isomorfismo local no implica que
lo mismo ocurra con~$f$. En efecto, si~$f$ es étale, $f^\an$ es étale y, por
tanto, un isomorfismo local~\hbox{\cite[\no13 \S 1]{XII.4}}, pero~$f$ no lo es
necesariamente.
\item[b)]
\ifthenelse{\boolean{orig}}
{el enunciado~\Ref{XII.3.2}}
{El enunciado~\Ref{XII.3.2}}%
~no es cierto si no se supone que~$f$ es de tipo finito. Demostremos, por
ejemplo, que~$f^\an$ puede ser una inmersión cerrada sin que~$f$ lo sea. Basta
tomar como~$X$ la suma disjunta de~$\ZZ$ copias de~$\Spec\CC$, como~$Y$ la
recta afín y como~$f$ el morfismo obtenido enviando los puntos de~$X$ a
puntos distintos de~$Y$ que forman un subconjunto discreto.
\end{enumerate}
\end{remarque}
